\begin{abstract}
Inference leakage depends on the background information available to an adversary. A field with little predictive value alone may substantially improve inference of a sensitive target when combined with other released fields. We propose an audit that turns set-level inferability into a field-level risk profile under an explicit background budget. Budgeted marginal risk is the largest gain over admissible backgrounds relative to already-public data; background amplification isolates the part missed by single-field assessment. We establish the measure's safety-margin interpretation and its connection to minimal unsafe combinations. To evaluate the many required field sets, an amortized estimator combines reconstruction-trained and random features with a separate ridge readout for each set. The scan proposes backgrounds, and dedicated retraining verifies selected gains and criticality evidence. Evaluation across four power-system data sets, ten targets and 124,425 exhaustively retrained field sets shows that average contextual gain is typically much smaller than its maximum. The estimator achieves a mean absolute marginal-risk error of 0.015 and a field-ranking correlation of 0.968. Verifying three backgrounds per field recovers 97.8\% of total marginal risk; critical-field recall is 76--91\% without false positives in the evaluated settings. At 48--64\,ms per set, the estimator recovers comparable certified marginal risk to TabPFN v2 at 13--110 times lower scan cost, while their threshold-specific recall differs. A controlled dispatch case and a withholding application demonstrate how the profile explains and supports decisions about data release.
\end{abstract}
